% Owners: both workload pairs. B selection is not inferred from these results.
\section{Conclusion, Lessons Learned, and Difficulties}
\label{sec:conclusion}

The project demonstrated the importance of identifying where transactions
spend time and aligning data placement with their access patterns.

For A, writes formed 75\% of the mix, but T3 and T4 accounted for approximately
94--97\% of the preserved baseline's recorded time. Our final design, a112,
combines lookup indexes, join-key colocation and read-only declarations.
Its matched seed-42 experiment achieved 54.238--90.788 TPS versus
3.225--3.613 TPS for the fresh baseline: an observed $16.8$--$28.2\times$
improvement, with zero failed measured a112 transactions.

The seed-7 repeat broadly reproduced throughput, but its aborted 64-client
warmup shows that concurrency reliability remains unresolved. Faster a115
observations also came with repeated aborts and no complete successful matrix.
We selected a112 while retaining these limitations.

Kelvin's write-path experiments distinguish mechanism improvements from
workload improvements. Autocommit reduced targeted-update median latency by
29--74\%, and fillfactor 80 produced 100\% HOT updates. Neither removed the
expensive reads on the a102 base. A10 remains a separate measured result;
no a10/a112 combination has been benchmarked.

For B, matched b0/b6 jobs achieved 108.250 versus 5.190 TPS at one client,
an observed $20.86\times$ improvement. At higher concurrency, aborted phases
and baseline failures limited comparisons. Separate correctness testing
provided evidence for inventory, order totals, identifiers and T7 conservation,
but did not retrospectively validate every earlier benchmark.

Three lessons stand out. First, optimization moves the bottleneck: once T3/T4
became cheaper, T6 dominated A. Second, added complexity is not automatically
beneficial: a114's maintained summaries did not improve its diagnostic,
while a115's simpler timestamp index was faster in successful cases.
Third, controls materially affect interpretation. The fresh and preserved
baselines differed substantially, so single runs on different allocations
cannot establish precise component-level gains.

Practical difficulties included distinguishing Slurm completion from pgbench
success, diagnosing worker errors surfaced as Citus shard-commit failures,
and maintaining comparable conditions. We examined phase exits, transaction
logs, failure counts and validation together. Future work should address
recurring T1 commit failures, repeat matched experiments, and test selected
write-path improvements on a112. Those combinations remain proposals.
